\documentclass[11pt,a4paper]{article}

% ----- Encoding & fonts (clean sans-serif) -----
\usepackage[utf8]{inputenc}
\usepackage[T1]{fontenc}
\usepackage{helvet}
\renewcommand{\familydefault}{\sfdefault}
\usepackage{microtype}

% ----- Math + formatting -----
\usepackage{amsmath,amssymb}
\usepackage{enumitem}
\usepackage{array}
\usepackage{tabularx}
\usepackage{booktabs}

% ----- Figures + drawing -----
\usepackage{graphicx}
\usepackage{float}

% ----- Code listings -----
\usepackage{listings}
\usepackage[table]{xcolor}

% ----- Layout -----
\usepackage[margin=1in,headheight=14pt]{geometry}
\usepackage{parskip}
\usepackage{titlesec}
\usepackage{fancyhdr}
\usepackage[hidelinks]{hyperref}

% ----- Theme colours -----
\definecolor{accent}{HTML}{1F4E79}     % deep professional blue
\definecolor{accentdk}{HTML}{2E2E2E}   % near-black for subsections
\definecolor{accentlt}{HTML}{D9E2F0}   % light blue table tint
\definecolor{rulegray}{HTML}{B9B9B9}
\definecolor{codebg}{HTML}{F7F7F4}
\definecolor{codecomment}{HTML}{6A737D}
\definecolor{codekw}{HTML}{0B5FA5}
\definecolor{codestr}{HTML}{B5210F}

% ----- Section / subsection styling -----
\titleformat{\section}
  {\color{accent}\Large\bfseries}{}{0pt}{}
  [\vspace{2pt}{\color{accent}\titlerule[1.1pt]}]
\titlespacing*{\section}{0pt}{16pt}{8pt}

\titleformat{\subsection}
  {\color{accentdk}\large\bfseries}{}{0pt}{}
\titlespacing*{\subsection}{0pt}{12pt}{3pt}

% ----- Page header / footer -----
\pagestyle{fancy}
\fancyhf{}
\fancyhead[L]{\small\color{accentdk}Exp. FPGA 5 -- Keyboard}
\fancyhead[R]{\small\color{accentdk}FPGA Lab}
\fancyfoot[C]{\small\thepage}
\renewcommand{\headrulewidth}{0.4pt}
\renewcommand{\footrulewidth}{0pt}

% ----- Verilog listing style -----
\lstdefinestyle{vlog}{
    language=Verilog,
    backgroundcolor=\color{codebg},
    basicstyle=\ttfamily\footnotesize,
    keywordstyle=\color{codekw}\bfseries,
    commentstyle=\color{codecomment}\itshape,
    stringstyle=\color{codestr},
    numbers=left,
    numberstyle=\tiny\color{gray},
    numbersep=8pt,
    frame=single,
    rulecolor=\color{rulegray},
    breaklines=true,
    showstringspaces=false,
    tabsize=4,
    xleftmargin=2.2em,
    framexleftmargin=2.2em,
    aboveskip=10pt,
    belowskip=6pt,
    captionpos=b,
}
\lstset{style=vlog}

% ----- A placeholder box for screenshots still to be pasted in -----
% Replace each \placeholder{...} with a real \includegraphics{...} once the
% image is exported from Vivado / taken with a camera.
\newcommand{\placeholder}[2][0.92\linewidth]{%
  \begin{center}
  \setlength{\fboxsep}{0pt}%
  \fcolorbox{rulegray}{white}{\begin{minipage}[c][4.4cm][c]{#1}
    \centering\itshape\color{gray}
    [ paste screenshot here ]\\[5pt]
    \footnotesize #2
  \end{minipage}}
  \end{center}%
}

% ----- A real figure: framed image + italic caption underneath -----
\newcommand{\reportfig}[3][\linewidth]{%
  \begin{center}
  \setlength{\fboxsep}{1.5pt}%
  \fcolorbox{rulegray}{white}{\includegraphics[width=#1]{#2}}\\[4pt]
  {\footnotesize\itshape #3}
  \end{center}%
}

\setlength{\parindent}{0pt}
\setlength{\parskip}{0.6em}

\begin{document}

% ======================================================================
%  TITLE BLOCK
% ======================================================================
\thispagestyle{fancy}
\begin{center}
  {\color{accent}\rule{\linewidth}{1.4pt}}\\[10pt]
  {\fontsize{24}{28}\selectfont \textbf{\color{accent}FPGA Experiment Report}}\\[6pt]
  {\large Experiment 5 --- PS/2 Keyboard Interface}\\[8pt]
  {\color{accent}\rule{\linewidth}{1.4pt}}
\end{center}

\vspace{6pt}

\renewcommand{\arraystretch}{1.3}
\begin{center}
\begin{tabular}{@{}l l l@{}}
\textbf{Task:} & FPGA 5 --- Keyboard (PS/2 interface) & \\
\textbf{Student 1:} & Mahmood Stitia & \textbf{ID:} 214032682 \\
\textbf{Student 2:} & Saleh Khalil  & \textbf{ID:} 213310485 \\
\end{tabular}
\end{center}

\vspace{6pt}

% ======================================================================
%  PART 1 - STATUS UPDATE
% ======================================================================
\section*{Part 1 --- Status update}

The table below lists the submitted files and their status. The status is one of
\emph{Done} (module and simulation written and passing), \emph{Partially done}
(started but not fully working --- explained in the notes), or \emph{Not started}.

\vspace{4pt}
\renewcommand{\arraystretch}{1.5}
\begin{tabularx}{\textwidth}{|>{\bfseries\ttfamily}l|c|X|}
\hline
\rowcolor{accent}
\textcolor{white}{\textbf{File / Step}} & \textcolor{white}{\textbf{Status}} & \textcolor{white}{\textbf{Notes / comments}} \\ \hline
\rowcolor{accentlt}
Ps2\_Interface.v     & Done & PS/2 receiver; testbench passes \\ \hline
Ps2\_Interface\_tb.v & Done & Self-checking; all scenarios pass \\ \hline
\rowcolor{accentlt}
Ps2\_Display.v       & Done & 7-seg + LED strobe (CDC to 100\,MHz) \\ \hline
Ps2\_Top.v           & Done & Top level, FPGA pin wiring \\ \hline
\rowcolor{accentlt}
task1.xdc            & Done & Pin + clock constraints \\ \hline
\end{tabularx}
\renewcommand{\arraystretch}{1.0}

\subsection*{Issues}

Any errors present at the time of submission are described here (per module /
simulation, with what the error means).

\begin{itemize}[topsep=2pt,itemsep=2pt,leftmargin=1.4em]
  \item \texttt{Ps2\_Interface.v}:
  \item \texttt{Ps2\_Interface\_tb.v}:
  \item \texttt{Ps2\_Display.v}:
  \item \texttt{Ps2\_Top.v}:
  \item \texttt{task1.xdc}:
\end{itemize}

% ======================================================================
%  PART 2 - SIMULATIONS AND ANSWERS
% ======================================================================
\section*{Part 2 --- Simulations and answers}

This part contains the simulation of \texttt{Ps2\_Interface} (annotated
waveforms) followed by the answers to the lab's questions. \emph{(Per the PDF,
only \texttt{Ps2\_Interface.v} is simulated.)}

\subsection*{Simulation --- Ps2\_Interface}

\placeholder[\linewidth]{Annotated waveform of \texttt{Ps2\_Interface\_tb}:
PS2Clk, rstn, PS2Data, scancode[7:0], keyPressed. Draw arrows on: the reset test,
the first make-code, the typematic repeats, the extended key, and the
release / re-press.}

\textbf{What the waveform shows} (annotate the real screenshot with arrows):
\begin{enumerate}[noitemsep,topsep=2pt,label=\arabic*.,leftmargin=1.6em]
  \item \textbf{Reset test} --- with \texttt{rstn=0} the outputs are cleared even while \texttt{PS2Clk} is idle (the async-reset advantage).
  \item \textbf{First make-code \texttt{0x2E}} --- after the stop-bit edge \texttt{scancode} becomes \texttt{2E} and \texttt{keyPressed} pulses exactly once.
  \item \textbf{Typematic repeats of \texttt{0x2E}} --- the same code is sent twice more; \texttt{scancode} stays \texttt{2E} and \texttt{keyPressed} does \emph{not} pulse again.
  \item \textbf{Extended key \texttt{0xE0,0x5A}} --- the \texttt{0xE0} frame leaves \texttt{scancode} unchanged and produces no pulse; the trailing \texttt{0x5A} sets \texttt{scancode=5A} with one pulse.
  \item \textbf{Release \texttt{0xF0,0x5A}} --- neither frame pulses; the displayed code is held.
  \item \textbf{Re-press \texttt{0x5A} after release} --- pulses again (a new press after a release is a new event).
\end{enumerate}

\placeholder[0.7\linewidth]{Zoom-in showing that the data is sampled on the
\emph{negative edge} of PS2Clk (arrow on a falling edge where scancode / keyPressed
react).}

\subsection*{Answer --- Asynchronous reset (advantages \& disadvantages)}

A flip-flop can be reset \emph{synchronously} (the reset is sampled only on the
active clock edge) or \emph{asynchronously} (the reset appears in the sensitivity
list and acts the instant it is asserted):
\begin{center}
\begin{tabular}{ll}
\texttt{always @(posedge clk)} & {\footnotesize synchronous: reset acts only on a clock edge} \\
\texttt{always @(posedge clk or negedge rstn)} & {\footnotesize asynchronous: reset acts immediately} \\
\end{tabular}
\end{center}

\renewcommand{\arraystretch}{1.3}
\begin{tabularx}{\textwidth}{|>{\bfseries}p{2.7cm}|X|X|}
\hline
\rowcolor{accent}
 & \textcolor{white}{\textbf{Asynchronous reset}} & \textcolor{white}{\textbf{Synchronous reset}} \\ \hline
Advantage & Resets the logic \emph{immediately}, even when \emph{no clock edge is available}. Critical in this lab because \texttt{PS2Clk} is idle most of the time. & Easy static-timing analysis; the reset is ``just another data signal'', no special routing, no recovery / removal checks. \\ \hline
Disadvantage & Reset \emph{removal} (de-assertion) is dangerous: releasing it too close to a clock edge violates recovery / removal time and can cause \textbf{metastability}; needs a reset synchroniser and dedicated async-reset routing. & Needs a toggling clock to take effect --- useless while the clock is idle; the reset pulse must be at least one clock period wide to be captured. \\ \hline
\end{tabularx}
\renewcommand{\arraystretch}{1.0}

\textbf{Why it matters here.} The keyboard clock \texttt{PS2Clk} toggles only
during a transmission and is idle (high) otherwise. A synchronous reset on that
domain would not act while the clock is idle, so \texttt{Ps2\_Interface} uses an
\textbf{active-low asynchronous reset} \texttt{rstn} (see Figure~1 of the
assignment: the push-button \texttt{reset} is inverted to
\texttt{\textasciitilde reset}/\texttt{rstn}). The de-assertion hazard is handled
on the fast side by a multi-FF synchroniser (see \texttt{Ps2\_Display}).

\subsection*{Answer --- PS/2 interface (frame, non-toggling clock, key-press detection)}

The PS/2 keyboard is a device-driven serial interface on two wires:
\texttt{PS2Clk} (clock generated \emph{by the keyboard}, $\approx$10--17\,kHz) and
\texttt{PS2Data}. One byte is sent as an \textbf{11-bit frame, LSB first}, and the
host samples \texttt{PS2Data} \textbf{on the falling edge of} \texttt{PS2Clk}:

\begin{center}
\renewcommand{\arraystretch}{1.3}
\begin{tabular}{|c|c|c|c|}
\hline
\rowcolor{accentlt}
\textbf{Bit 0} & \textbf{Bits 1--8} & \textbf{Bit 9} & \textbf{Bit 10} \\ \hline
Start (0) & Data D0..D7 (LSB first) & Parity (odd) & Stop (1) \\ \hline
\end{tabular}
\renewcommand{\arraystretch}{1.0}
\end{center}

Key facts used by the design:
\begin{itemize}[noitemsep,topsep=2pt,leftmargin=1.4em]
  \item \textbf{Make code} --- sent on key press (e.g.\ keypad \texttt{5} $\rightarrow$ \texttt{0x2E}).
  \item \textbf{Typematic repeat} --- while a key is held, the make code is re-sent repeatedly; \texttt{keyPressed} must pulse \emph{only once} (on the first press).
  \item \textbf{Break code} --- on release: \texttt{0xF0} followed by the make code.
  \item \textbf{Extended keys} --- prefixed with \texttt{0xE0}; the \texttt{0xE0} must be ignored and only the trailing byte displayed.
\end{itemize}

\textbf{Non-toggling clock $\Rightarrow$ decode early.} Because \texttt{PS2Clk}
stops after the stop bit, the module must \emph{not} wait for an edge after the
frame. \texttt{Ps2\_Interface} decodes the byte \emph{on} the stop-bit edge (the
last edge the keyboard actually produces), using a 22-bit shift register that
holds the current frame and the previous one so two-frame sequences
(\texttt{0xE0}/\texttt{0xF0}) can be recognised.

\textbf{Two clock domains (CDC).} \texttt{scancode}/\texttt{keyPressed} are
produced in the slow \texttt{PS2Clk} domain and consumed in the fast 100\,MHz
\texttt{clk} domain of \texttt{Ps2\_Display}. \texttt{keyPressed} is passed
through a 3-FF synchroniser and edge-detected; \texttt{scancode} is captured on
that already-settled pulse --- a simple, safe data + valid CDC handshake. This is
also why the logic analyser cannot debug both domains in a single ILA core.

% ======================================================================
%  ELABORATED DESIGN
% ======================================================================
\section*{Elaborated Design}

\subsection*{Top module}

\reportfig[\linewidth]{linter_elaborated_design.png}{RTL-elaborated schematic of
\texttt{Ps2\_Top}: \texttt{u\_interface} (\texttt{Ps2\_Interface}, PS2Clk domain)
feeds \texttt{keyPressed} and \texttt{scancode[7:0]} into \texttt{u\_display}
(\texttt{Ps2\_Display}, 100\,MHz domain); the \texttt{reset} push-button passes
through an \texttt{RTL\_INV} (\texttt{rstn\_i}) into \texttt{rstn}. Inputs
\texttt{clk, PS2Clk, PS2Data, reset} are on the left and outputs
\texttt{an[3:0], dp, led, seg[6:0]} on the right.}

\subsection*{Ps2\_Interface}

\reportfig[\linewidth]{sch_innterface.png}{RTL-elaborated schematic of
\texttt{Ps2\_Interface}: the 22-bit shift register and the \texttt{bitcount}
counter feed the \texttt{held\_code} / \texttt{key\_down} registers and the
make / break / extend decode logic, all clocked by \texttt{PS2Clk}; the outputs
\texttt{scancode[7:0]} and \texttt{keyPressed} leave on the right.}

% ======================================================================
%  IMPLEMENTATION
% ======================================================================
\section*{Implementation}

\subsection*{Synthesis schematic}

\reportfig[\linewidth]{synthesis_schematic.png}{Post-synthesis schematic. Every
external port now passes through an \texttt{IBUF}/\texttt{OBUF}, the two clocks
(\texttt{clk}, \texttt{PS2Clk}) go through \texttt{BUFG} global buffers, and
\texttt{u\_interface} / \texttt{u\_display} drive the \texttt{an}, \texttt{dp},
\texttt{led} and \texttt{seg} output buffers.}

\subsection*{Project summary}

\reportfig[\linewidth]{project_summary.png}{Post-implementation project summary.
No DRC violations; timing met (WNS\,$=5.269$\,ns, 0 failing endpoints);
utilization is small (LUT \& FF $\approx$1\%, IO 16\%, BUFG 6\%); total on-chip
power is $0.076$\,W with a junction temperature of $25.4^{\circ}$C.}

\subsection*{Timing summary}

\reportfig[\linewidth]{design_timing_report.png}{Design Timing Summary after
implementation. The Setup (WNS\,$=5.269$\,ns), Hold (WHS\,$=0.330$\,ns) and
Pulse-Width (WPWS\,$=4.500$\,ns) slacks are all positive with 0 failing
endpoints --- \emph{``All user specified timing constraints are met''}.}

\subsection*{Critical Path}

\reportfig[\linewidth]{critical_path.png}{Intra-Clock Paths report. The worst
paths (Path~1--4) run inside \texttt{u\_display} from
\texttt{led\_timer\_reg[16]/C} to \texttt{led\_timer\_reg[20..23]/CE}, with slack
$5.269$\,ns, 3 logic levels and $4.361$\,ns total delay ($1.054$\,ns logic
$+\,3.307$\,ns net) against the $10$\,ns requirement --- comfortably met.}

\subsection*{Device}

\reportfig[0.55\linewidth]{synthesized_design.png}{Device (floorplan) view of the
placed-and-routed design on the Artix-7 (\texttt{XC7A35T}). The few used slices
are highlighted across the clock regions (\texttt{X0Y0}--\texttt{X1Y2}); the
design occupies only a tiny fraction of the fabric.}

% ======================================================================
%  PROBLEMS WE FACED
% ======================================================================
\section*{Problems we faced}

On hardware the design behaves as specified: pressing a numeric-keypad key shows
its 8-bit make-code as two hex symbols on the two right-hand 7-segment digits
(e.g.\ \texttt{5}\,$\rightarrow$\,\texttt{2E}); each new press blinks the LED once
with a short strobe (holding the key does \emph{not} re-blink); the displayed code
latches until the next press; and extended keys show only their last two hex
symbols (the \texttt{0xE0} prefix is invisible).

\emph{(Describe the problems you actually hit in the lab and how you solved them.
For example:)}
\begin{itemize}[noitemsep,topsep=2pt,leftmargin=1.4em]
  \item \emph{(fill in)} where to put the \texttt{0xF0}/\texttt{0xE0} / two-key handling logic --- in the interface vs.\ the display module --- and the wrong display behaviour seen before moving it all into \texttt{Ps2\_Interface}.
  \item \emph{(fill in)} the keyboard clock being idle most of the time (decode before the stop bit, async reset).
  \item \emph{(fill in)} CDC glitches on \texttt{keyPressed} and debugging the two clock domains in separate ILA cores.
\end{itemize}

\placeholder[0.7\linewidth]{Board photo: keypad plugged into the USB host port, a
key pressed and its hex code shown on the two right digits with the LED lit.}

% ======================================================================
%  APPENDIX - VERILOG SOURCE CODE
% ======================================================================
\section*{Appendix --- Verilog Source Code}

The full source of the submitted files is listed below. Paste the code of each
file inside its block.

\subsection*{Ps2\_Interface.v}
\begin{lstlisting}
// --- paste Ps2_Interface.v here ---
\end{lstlisting}

\newpage
\subsection*{Ps2\_Interface\_tb.v}
\begin{lstlisting}
// --- paste Ps2_Interface_tb.v here ---
\end{lstlisting}

\newpage
\subsection*{Ps2\_Display.v}
\begin{lstlisting}
// --- paste Ps2_Display.v here ---
\end{lstlisting}

\newpage
\subsection*{Seg\_7\_Display.v}
\begin{lstlisting}
// --- paste Seg_7_Display.v here ---
\end{lstlisting}

\newpage
\subsection*{Ps2\_Top.v}
\begin{lstlisting}
// --- paste Ps2_Top.v here ---
\end{lstlisting}

\newpage
\subsection*{task1.xdc}
\begin{lstlisting}[language=]
## --- paste task1.xdc here ---
\end{lstlisting}

\end{document}
